\documentclass[twocolumn, a4j]{jarticle}

% 画像の挿入関連のパッケージ
\usepackage[dvipdfmx]{graphicx}
\usepackage[dvipdfmx]{color}

% 数式関連のパッケージ
\usepackage{amsmath}
\usepackage{amssymb}
\usepackage{array}

% URLをきれいに入れるためのパッケージ
\usepackage{url}

% 松岡研のパッケージ（ヘッダータイトル関連のパッケージ）
\usepackage{M-lab}

% その他，好きなパッケージを入れることができる
\usepackage{balance}

% 自分でコマンド書いたり改造したりもできる
\newcommand{\redb}[1]{\textbf{\textcolor{red}{#1}}}

%%%%%%%%%%%%%%% 以下，良しなに編集してください． %%%%%%%%%%%%%%%%%%%
\Author{瀧口 昇汰}
\aAuthor{Shouta Takiguchi}
\ID{C0B23092}
\Title{物体追跡研究用データセット作成シミュレータの開発}
\Laboratory{松岡}
\HeaderSubject{2026年度 コンピュータサイエンス学部 卒業課題I 中間発表}
% \HeaderSubject{2025年度 コンピュータサイエンス学部 卒業課題II 最終発表}
\Date{2026年08月03日}

\begin{document}
\twocolumn[\makeHeader]

%%%%%%%%%%%%%%% 本文 %%%%%%%%%%%%%%%%%%%
\section{背景}
物体トラッキング技術の性能を評価するためには，トラッキング対象の正確な位置および姿勢を正解値として付与したテストデータが多く必要となる．特に，近年の手法では高精度な評価が求められるため，大量かつ高品質なデータセットの構築が重要な課題となっている．

しかし，現状の画像または映像からテストデータの正解値を設定する場合，移動する物体の位置や姿勢を連続的にかつ高精度に設定することは容易ではない．
例えば，人物が他の物体の背後に隠れる遮蔽や，屋内外の移動による急激な照明変動が発生する環境では，対象の外観が影や反射によって大きく変化する．このような状況では自動的に正確な位置や姿勢を推定することが困難となるため，手動で正解値を設定する必要がある．その結果，データ数が増加するほど作業量が増大し，膨大な時間と労力を要するという問題が生じる．

\section{既存技術と課題}
このような背景から，物体トラッキングの評価時にはシミュレータを用いて生成したテストデータを利用する手法が広く採用されている．シミュレータを使用することで，物体の座標や向きといった情報を正確に取得できるため，正解値付きの画像系列を効率的に生成することが可能となる．

しかし，既存のシミュレータを用いてテストデータを作成するためには，トラッキング対象の移動経路を事前に設計する必要がある．この設計作業にはシミュレーション環境の構築やプログラミングに関する専門的な知識が求められ，その分野に精通していない研究者にとっては大きな負担となる．特に，複雑な動きや多数の対象を扱う場合には，設計の難易度がさらに増加する．

さらに，シミュレータで生成された画像と実環境で撮影された画像との間に写実的な差異が大きい場合，シミュレーション上で得られた評価結果が実環境において適切に再現されない可能性がある．そのため，背景や照明条件を含む環境全体とトラッキング対象を，実環境に近い外観で再現することが重要である．

人流および歩行者シミュレーションに関する既存研究・既存システムとして，Oasys MassMotionやPTV Viswalkが挙げられる．

Oasys MassMotionは，群衆流動の解析や設計段階での安全性・効率性の検討を目的としたシミュレーションソフトウェアであり，複数のシナリオ比較や3次元環境における可視化に対応している．また，2Dおよび3D形状の利用や，各種CAD/BIMデータの取り込みに対応しており，大規模空間における人流解析に活用されている．
一方で，Oasys MassMotionはあくまでも群衆流動の解析に特化しており，フォトリアリスティック\footnote{フォトリアリスティックとは，人間が実写と見分けがつかないほど写実的に表現する技術やその状態を指す．}に表現した画像を出力することはできない．

PTV Viswalk は，歩行者一人一人の挙動を扱うシミュレーションソフトウェアであり，歩行者同士の相互作用に加えて，自動車，バス，鉄道，自転車など他の交通主体との関係も含めて解析できる．これにより，通路幅，待機領域の大きさ，移動時間，待ち時間，混雑状況，安全性などの観点から評価を行うことが可能であり，駅，空港，都市空間，イベント会場などにおける人流設計に利用されている．
しかし，本格的なシミュレーションを行うためには，プログラミングを含む専門的な知識が多く求められるという課題がある．

\section{研究目的}
本研究では，上記の課題を踏まえ，物体トラッキング手法の評価に使用可能な画像系列および正解値を容易に生成できるテストデータ生成シミュレータの開発を目的とする．本システムは以下の要件を満たすことを目指す．

一つ目は，専門的なプログラミング知識を必要とせずにトラッキング対象の動きを設計できることである．これにより，非専門家でも直感的にデータ生成を行える環境を提供する．

二つ目は，フォトリアリスティックに環境を表現することである．これにより，実環境との乖離を最小限に抑えることで評価の信頼性向上を図る．

また，研究室内で実施されている物体トラッキング研究の経験に基づき，目的地まで並ぶ対象となる列の簡易な構築機能や，列を横切るノイズとなる動物体の生成機能を実装する．
これにより現実に近い複雑な状況を再現可能とし，より実用的な評価環境の構築を目指す．最終的には，テストデータ作成の負担を軽減するとともに，実環境に近い条件での性能評価を可能とすることを目的とする．

\section{研究のアプローチ}
\begin{figure}[tb]
    \centering \includegraphics[width=0.5\textwidth]{fig/godot.png}
    \vspace{-1em}
    \caption{Godotのドロワー画面}
    \label{fig:Main.tscnのメイン画面}
\end{figure}

\begin{figure}[tb]
    \centering \includegraphics[width=0.5\textwidth]{fig/kari.png}
    \vspace{-1em}
    \caption{3D変換処理後の画像}
    \label{fig:3D変換処理後の画像}
\end{figure}

本研究で開発するシステムは，最終的に三つの手順で動作することを想定している．

まず，使用者は動物体の速度や加速度などのパラメータを数値として入力する．次に，ドロワー機能を用いて2次元空間上に線を描画し，動物体の移動経路を直感的に設計する．その後，入力された情報と経路データを基に軽量な移動シミュレーションを実行し，各時刻における位置および向きを算出したデータを出力する．最後に，得られた時系列データを3次元空間上に変換し，フォトリアリスティックな映像として出力する．

この一連の処理において，動物体の動きを誰でも容易に設計できることが重要であり，本研究の本質的な目的となる．

\section{評価項目}
本システムの評価は，松岡研で開発されている「列検知型人流計測システム」を用いて行う．評価指標としては，実際の動画を使用した場合と同等の性能指標が得られることを目標とする．

\section{これまでに行ったこと}
これまでに，ソフトウェア全体の構造設計を行い，Godot 4.7を用いたプロトタイプの開発を進めている．

ドロワーによる移動経路設計において，意図した動きを再現できない問題が確認された．また、経路編集時の操作性にも課題があり，ユーザビリティの観点からの改善が必要となる．現在はこれらの問題を解決するための別手法について検討を行っている．

\section{今後の研究計画}
今後の研究計画として，11月までに本システムをプロトタイプからベータ版として完成させることを目標とする．その過程で，認識システムを用いた評価を繰り返し行い，認識率の向上を確認しながら改良を加えていく．さらに，得られた結果を基にシステムの有効性を検証し，卒業論文としてまとめる予定である．

\begin{center}
    \refstepcounter{table}
    \label{tab:schedule}
    表\thetable：今後の研究計画
    \begin{tabular}{|c|p{50mm}|}
        \hline
        時期 & 実施内容 \\
        \hline
        8月 9月 & プロトタイプの完成 \\
        \hline
        10月 & ベータ版の完成 \\
        \hline
        11月 & 評価項目の解決 \\
        \hline
        12月 1月 & 卒業論文の執筆 \\
        \hline
    \end{tabular}
\end{center}

% 参考文献
% mybib.bibを参考に参考文献を追加すること。
% 1つも参照されていない場合エラーが出るので、自分で参考文献を追加するまではテンプレート内の\citeを消さないことを推奨する。
\bibliographystyle{junsrt}
\bibliography{mybib}
\end{document}

